\section{Comparison with Three-Layered Architecture}
\label{sec:comparison-with-three-layered-architecture}

The principles in Section~\ref{sec:key-principles-of-local-first-software} can be understood through a familiar application structure: presentation, business logic, and data access.
These responsibilities remain useful in local-first software.
What changes is where ordinary reads and writes happen, which copy of the data they use, and how changes reach other users.
This chapter follows a shared task-board example to explain these differences.

\subsection{The familiar three-layered model}

In \emph{Patterns of Enterprise Application Architecture}, Fowler describes three principal layers: presentation, domain, and data source.
The presentation layer displays information and handles user interaction.
The business-logic layer, also called the domain layer, implements application rules and calculations.
The data-source layer communicates with databases and other systems that provide services to the application.
For the task-board example, its main responsibility is data access: reading and writing persistent tasks.
\parencite[19--20]{patterns-of-enterprise-application-architecture}

Separating these responsibilities allows developers to understand one layer without knowing the details of the others.
Fowler requires the domain and data-source layers to remain independent of presentation, but permits presentation to access the data source directly.
The dependency between domain logic and the data source depends on the patterns chosen for data access.
The three layers therefore describe a separation of responsibilities without requiring every operation to pass through a strict chain of three modules.
\parencite[17, 20--22]{patterns-of-enterprise-application-architecture}

A layer describes a responsibility in the software; a tier describes where software runs.
Three logical layers do not require three separate machines, and all three can run on one device.
Fowler explicitly gives the example of running all three layers on a laptop with a local database.
\parencite[19]{patterns-of-enterprise-application-architecture}
Layering alone therefore says nothing about whether the application works offline or treats a server as authoritative.

For the comparison, consider a common server-centred deployment of these layers.
The browser displays a task board, an application server applies business rules, and a database stores tasks.
The server checks permissions and validates requested changes before saving them through the data-access layer.
The browser may keep recently fetched tasks in memory, but the server's database determines the accepted shared state.
This is the cloud data model described by \textcite{local-first-software}.
The comparison below uses this deployment as its baseline, rather than assuming that all layered applications require a server.

\subsection{Moving ordinary work to the device}

Fowler already identifies responsiveness and disconnected operation as reasons to run logic on the client.
He also notes that disconnected use can require a local data source whose changes must later be synchronized with the server, while leaving the synchronization problem outside the book's scope.
This provides a useful starting point for understanding local-first software: the familiar layers can move onto the device, but the application must also support independently editable copies and their reconciliation.
\parencite[23--24]{patterns-of-enterprise-application-architecture}

In a local-first version of the task board, each participating device holds an independently editable copy of the relevant board data.
Such a copy is called a replica.
The browser reads tasks from its local replica and applies edits locally, without waiting for a remote request.
Persistent storage keeps those edits available across application restarts.
A synchronization mechanism exchanges changes with other replicas when a connection is available.
\parencite{local-first-software}

In \emph{Designing Data-Intensive Applications}, \textcite[Chapter~5]{designing-data-intensive-applications} describes offline clients as a use case for multi-leader replication: several replicas can accept writes independently and exchange them asynchronously.
Each device's local database acts as a leader for its own writes, and synchronization may be delayed for hours or days while the device is disconnected.
He also explains collaborative editing as a replication problem, with edits applied to a local document replica and then shared with other participants.

This gives a distributed-systems explanation of how the local-first task board can support independent editing across devices.
The ownership and long-term access principles discussed earlier remain additional requirements beyond the replication mechanism.

The familiar responsibilities still have a place in this design.
Presentation renders the board and handles input.
Business logic checks rules that can be evaluated from local information, such as requiring a non-empty task title.
Data access loads and saves the local data.
These responsibilities can remain separate modules even though they now execute on the same device.
The client must contain enough logic and data to support the operations offered offline.

Synchronization adds a responsibility beyond local persistence.
Saving a task on a laptop does not, by itself, make that task available on a phone.
The application must track and exchange changes, apply incoming updates, and reconcile edits made independently.
Local persistence and synchronization have different failure conditions: a successful local save can coexist with an unavailable synchronization service.
\parencite[Chapter~5]{designing-data-intensive-applications}

A server can still help devices discover one another, relay updates, or retain a replica for devices that connect at different times.
Its presence does not make the application server-centred.
The important distinction is whether ordinary work with locally available data requires the server to respond.
In the local-first model used here, the server supports collaboration while the device remains able to read, edit, and save its own copy independently.
\parencite{local-first-software}

Figure~\ref{fig:layer-deployment-comparison} shows how the same responsibilities can be distributed across the two deployments.
The arrows represent runtime communication rather than source-code dependencies.
In the local-first deployment, synchronization connects the replicas while ordinary reads and writes stay within each device.

\begin{illustration}[htbp]
	\centering
	\caption{Placement of layers and persistent data in the two task-board deployments}
	\label{fig:layer-deployment-comparison}
	\begingroup
	\setstretch{1}
	\begin{tikzpicture}[
	thesis,
	layer/.style={thesis/box, minimum width=2.65cm, minimum height=0.6cm},
	store/.style={thesis/store, minimum width=2.65cm, minimum height=0.8cm},
	local/.style={thesis/local, <->},
	network/.style={thesis/network, <->}
]
	\node[thesis/heading] at (3.45,0) {(a) Server-centred};
	\node[thesis/heading] at (11.1,0) {(b) Local-first};

	\draw[thesis/boundary] (0.1,-0.5) rectangle (3.25,-1.95);
	\draw[thesis/boundary] (3.65,-0.5) rectangle (6.8,-1.95);
	\node at (1.675,-0.8) {Device A: browser};
	\node at (5.225,-0.8) {Device B: browser};
	\node[layer] (remote-ui-a) at (1.675,-1.4) {Presentation};
	\node[layer] (remote-ui-b) at (5.225,-1.4) {Presentation};

	\draw[thesis/boundary] (1.35,-3.25) rectangle (5.55,-5.85);
	\node at (3.45,-3.55) {Application server};
	\node[layer] (remote-domain) at (3.45,-4.15) {Business logic};
	\node[layer] (remote-data) at (3.45,-5.2) {Data access};
	\draw[local] (remote-domain) -- (remote-data);
	\draw[network] (remote-ui-a.south) -- (1.675,-2.4) -- (3.45,-2.4) -- (3.45,-3.25);
	\draw[network, <-] (remote-ui-b.south) -- (5.225,-2.4) -- (3.45,-2.4);
	\node[thesis/label] at (3.45,-2.8) {Requests / responses};

	\draw[thesis/boundary] (1.35,-7.05) rectangle (5.55,-8.5);
	\node at (3.45,-7.35) {Database server};
	\node[store] (remote-store) at (3.45,-7.95) {Authoritative tasks};
	\draw[network] (remote-data.south) -- (3.45,-7.05);
	\node[thesis/label] at (3.45,-6.45) {Database reads / writes};

	\draw[thesis/boundary] (7.6,-0.5) rectangle (10.9,-5.6);
	\draw[thesis/boundary] (11.3,-0.5) rectangle (14.6,-5.6);
	\node at (9.25,-0.8) {Device A};
	\node at (12.95,-0.8) {Device B};
	\node[layer] (local-ui-a) at (9.25,-1.4) {Presentation};
	\node[layer] (local-domain-a) at (9.25,-2.45) {Business logic};
	\node[layer] (local-data-a) at (9.25,-3.5) {Data access};
	\node[store] (local-store-a) at (9.25,-4.85) {Persistent\\local replica};
	\node[layer] (local-ui-b) at (12.95,-1.4) {Presentation};
	\node[layer] (local-domain-b) at (12.95,-2.45) {Business logic};
	\node[layer] (local-data-b) at (12.95,-3.5) {Data access};
	\node[store] (local-store-b) at (12.95,-4.85) {Persistent\\local replica};
	\draw[local] (local-ui-a) -- (local-domain-a);
	\draw[local] (local-domain-a) -- (local-data-a);
	\draw[local] (local-data-a) -- (local-store-a);
	\draw[local] (local-ui-b) -- (local-domain-b);
	\draw[local] (local-domain-b) -- (local-data-b);
	\draw[local] (local-data-b) -- (local-store-b);

	\draw[thesis/boundary] (8.8,-7.05) rectangle (13.4,-8.5);
	\node at (11.1,-7.35) {Synchronization server};
	\node[store, minimum width=3.9cm] (sync-server) at (11.1,-7.95) {Relay and stored replica};
	\draw[network] (local-store-a.south) -- (9.25,-6.05) -- (11.1,-6.05) -- (11.1,-7.05);
	\draw[network, <-] (local-store-b.south) -- (12.95,-6.05) -- (11.1,-6.05);
	\node[thesis/label] at (11.1,-6.45) {Asynchronous synchronization};

	\draw[local] (0.35,-9.15) -- (1.15,-9.15);
	\node[anchor=west] at (1.25,-9.15) {Local interaction};
	\draw[network] (4.4,-9.15) -- (5.2,-9.15);
	\node[anchor=west] at (5.3,-9.15) {Network communication};
	\draw[thesis/boundary] (10.35,-9.3) rectangle (10.85,-9.0);
	\node[anchor=west] at (10.95,-9.15) {Deployment boundary};
\end{tikzpicture}

	\endgroup
	\figuresource{Author's illustration based on \textcite{patterns-of-enterprise-application-architecture,local-first-software,designing-data-intensive-applications}.}
\end{illustration}
\FloatBarrier

\subsection{Following a task update}

Suppose a user changes a task's title on a shared board.
In the server-centred application, the operation follows this sequence:

\begin{enumerate}
	\item The presentation layer sends the new title and task identifier to the application server.
	\item The server checks the user's permission and validates the title.
	\item The data-access layer writes the change to the database.
	\item The server returns a result, and the interface confirms that the change was saved.
\end{enumerate}

An optimistic interface can display the new title before the response arrives.
However, this displayed result remains provisional until the server accepts the request.
If the server rejects it, the interface must handle the rejection, for example by restoring the previous title or retaining the input for correction.
Hiding the wait does not remove the remote dependency.
\parencite{local-first-software}

In the local-first application, the corresponding operation is:

\begin{enumerate}
	\item The presentation layer passes the new title to local business logic.
	\item The application validates the title and updates its local replica.
	\item The interface reflects the change, and local persistence saves it on the device.
	\item Synchronization sends the change to other replicas when possible and incorporates their changes in return.
\end{enumerate}

The interface can update before persistence finishes, but it should report the edit as saved locally only after that save succeeds.
Once saved locally, the edit survives an ordinary application restart without requiring a successful upload.
The interface can distinguish this state from being synchronized with other devices.
Losing the connection delays sharing the title; it does not prevent saving it locally.
\parencite{local-first-software}

Figure~\ref{fig:task-update-sequences} separates the visible update from the successful save and subsequent synchronization.
Both examples begin with the same title edit; the local-first example assumes that local validation has passed.
In the illustrated synchronization protocol, the server acknowledges receipt of the change, allowing the client to update its synchronization status.
This acknowledgement does not imply that every other collaborator has already received the change.
The vertical order shows the sequence of events, not their duration.

\begin{illustration}[htbp]
	\centering
	\caption{Saving a title edit with server confirmation or local persistence}
	\label{fig:task-update-sequences}
	\begingroup
	\setstretch{1}
	\begin{tikzpicture}[
	thesis,
	participant/.style={thesis/participant, minimum height=0.65cm, text width=1.7cm},
	status/.style={thesis/box, text width=1.75cm, inner sep=3pt},
	lifeline/.style={thesis/lifeline},
	local/.style={thesis/local, ->},
	network/.style={thesis/network, ->},
	message/.style={thesis/label, above}
]
	\fill[thesis/region] (7.3,-1.35) rectangle (14.9,-5.75);
	\node[thesis/heading] at (3.6,0) {(a) Server-centred, optimistic display};
	\node[thesis/heading] at (11.1,0) {(b) Local-first, initially offline};
	\node[participant] at (1.1,-0.8) {Browser};
	\node[participant] at (3.65,-0.8) {Application\\server};
	\node[participant] at (6.2,-0.8) {Database\\server};
	\node[participant] at (8.35,-0.8) {Interface};
	\node[participant] at (11.1,-0.8) {Local data\\and sync};
	\node[participant] at (13.85,-0.8) {Sync\\server};

	\draw[lifeline] (1.1,-1.15) -- (1.1,-8.95);
	\draw[lifeline] (3.65,-1.15) -- (3.65,-8.95);
	\draw[lifeline] (6.2,-1.15) -- (6.2,-8.95);
	\draw[lifeline] (8.35,-1.15) -- (8.35,-8.95);
	\draw[lifeline] (11.1,-1.15) -- (11.1,-8.95);
	\draw[lifeline] (13.85,-1.15) -- (13.85,-8.95);

	\node[status] at (1.1,-1.9) {Display updated\\(provisional)};
	\draw[network] (1.1,-2.75) -- node[message] {Send title} (3.65,-2.75);
	\node[status] at (3.65,-3.35) {Validate};
	\draw[network] (3.65,-4.05) -- node[message] {Write title} (6.2,-4.05);
	\draw[network] (6.2,-4.75) -- node[message] {Persisted} (3.65,-4.75);
	\draw[network] (3.65,-5.45) -- node[message] {Confirm save} (1.1,-5.45);
	\node[status] at (1.1,-6.15) {Saved on server};
	\node[align=center, text width=6.3cm] at (3.6,-9.4) {Saving requires a successful network exchange.};

	\node[anchor=north east, inner sep=4pt] at (14.9,-1.35) {Disconnected};
	\node[status] at (8.35,-1.9) {Display updated};
	\draw[local] (8.35,-2.75) -- node[message] {Save title} (11.1,-2.75);
	\draw[local] (11.1,-3.65) -- node[message] {Persisted} (8.35,-3.65);
	\node[status] at (8.35,-4.35) {Saved locally};
	\node[align=center, text width=4.8cm] at (10.6,-5.15) {Editing can continue while disconnected.};
	\draw[thesis/divider] (7.3,-5.85) -- (14.9,-5.85);
	\node[thesis/label, inner sep=3pt] at (11.1,-5.85) {Connection restored};
	\draw[network] (11.1,-6.55) -- node[message] {Send changes} (13.85,-6.55);
	\draw[network] (13.85,-7.3) -- node[message] {Acknowledge} (11.1,-7.3);
	\draw[local] (11.1,-8.0) -- node[message] {Update sync status} (8.35,-8.0);
	\node[status] at (8.35,-8.65) {Synchronized\\with server};
	\node[align=center, text width=6.6cm] at (11.1,-9.4) {Local saving completes before sharing.};

	\draw[local] (1.5,-10.15) -- (2.3,-10.15);
	\node[anchor=west] at (2.4,-10.15) {Local interaction};
	\draw[network] (7.8,-10.15) -- (8.6,-10.15);
	\node[anchor=west] at (8.7,-10.15) {Network message};
\end{tikzpicture}

	\endgroup
	\figuresource{Author's illustration based on \textcite{local-first-software,designing-data-intensive-applications}.}
\end{illustration}

The same distinction applies to reads.
Reopening an already available board can load its local data without contacting the server.
Opening a board that has never been downloaded still requires obtaining that data, and offline use also requires the application's code to be available on the device.
Local-first operation removes the remote dependency for available work; it cannot supply information that the device has never received.

\subsection{Concurrent changes and business rules}

The server-centred model gives the application a common point at which to check and coordinate writes.
It still needs a concurrency policy: two users may submit changes based on the same earlier task title.
The server can, for example, use an atomic version check to reject an update based on an outdated task, or use locking to coordinate changes.
Layering itself does not choose this policy or prevent lost updates.
\parencite[Chapter~7]{designing-data-intensive-applications}

Local-first software must also handle changes made while replicas cannot communicate.
With asynchronous multi-leader replication, conflicting writes may both be accepted locally before either replica learns about the other write.
The replicas therefore need rules for reconciling those writes when they exchange changes.
\parencite[Chapter~5]{designing-data-intensive-applications}

Suppose one user renames a task on a disconnected laptop while another user moves the same task to a different column.
A merge mechanism that tracks the title and column independently can preserve both edits.
If both users change the title, the application needs a defined way to select or expose the competing values.
Replacing one entire board with the other would risk discarding unrelated work.

CRDTs provide one approach to this reconciliation.
They allow replicas that receive the same updates to converge, even when those updates arrive in different orders.
The JSON CRDT described by \textcite{json-crdt}, for example, supports concurrent changes to nested data structures.
Convergence does not mean that disconnected users always see identical state, nor does it guarantee that every merged result satisfies the application's business rules.

Business-rule violations can arise even when concurrent operations modify different records.
\textcite[Chapter~7]{designing-data-intensive-applications} illustrates this in his discussion of write skew: each transaction checks a condition and makes a change that appears valid, but their combined changes violate the condition.

An analogous problem arises when disconnected devices check a shared constraint independently.
Consider a board with a strict limit of five tasks in a column.
Two disconnected devices may each see four tasks and independently move one more task into that column.
Both local checks pass, but preserving both moves produces six tasks after synchronization.
The application must decide whether to allow a temporary excess, reconcile the moves later, or coordinate permission to use the remaining capacity.
If final approval requires contacting a server, that particular operation cannot be completed independently while offline.

\subsection{Responsibilities and trade-offs}

Table~\ref{tab:server-centred-and-local-first} summarizes the two task-board designs.
The differences concern data authority and communication, while the three logical responsibilities can remain present in both.

\begin{table}[htbp]
	\centering
	\caption{Comparison of the server-centred and local-first task boards}
	\label{tab:server-centred-and-local-first}
	\small
	\begin{tabular}{@{}>{\raggedright\arraybackslash}p{0.22\linewidth}>{\raggedright\arraybackslash}p{0.35\linewidth}>{\raggedright\arraybackslash}p{\dimexpr0.43\linewidth-4\tabcolsep\relax}@{}}
		\toprule
		Aspect         & Server-centred design                                      & Local-first design                                                             \\
		\midrule
		Working data   & Client data is subordinate to the server's accepted state. & Each device has an independently editable replica.                             \\
		Saving an edit & Completion depends on a server-side write.                 & Local saving completes independently of synchronization.                       \\
		Network outage & Remote operations must wait or remain provisional.         & Available local data remains readable and editable.                            \\
		Collaboration  & The server coordinates accepted updates.                   & Replicas exchange and reconcile independent updates.                           \\
		Business rules & The server can check shared state before accepting writes. & Offline checks use local knowledge; global constraints need a separate policy. \\
		Server role    & Required for authoritative reads and writes.               & Supports synchronization and storage of additional copies.                     \\
		\bottomrule
	\end{tabular}
\end{table}

The server-centred design concentrates control over shared state in the backend.
This is useful when an operation must be checked against current shared information before it can be accepted.
Local-first design gives the device more responsibility: it must preserve local work, handle incoming changes, and communicate the difference between local saving and synchronization.
Locally persisted edits survive application restarts, but can still be lost if the device's storage fails or is erased before they are synchronized or backed up elsewhere.
\parencite{local-first-software}

Three-layered architecture remains a useful way to organize either application.
To imagine the local-first version, place the presentation, the logic needed for ordinary editing, and persistent working data on each user's device, then connect those devices through synchronization.
The user works with a local replica, and communication brings the replicas together as connectivity permits.
This arrangement gives the preceding chapter's principles a concrete architectural form.
